% =============================================================================
%  Partie IV — Structures (Chapitre IV)
% =============================================================================
\section{Espèces de structures, isomorphismes, morphismes (IV)}

Le chapitre~IV introduit les \emph{espèces de structures} \textbf{sans nouvel
axiome} : ce sont des extensions définitionnelles de la théorie des ensembles.
La formalisation y certifie le \emph{squelette déductif} des critères de
transport de structure (CST1--CST23 de Bourbaki). Le choix d'ingénierie est
explicite et fidèle à la stratification : les schémas de transport (CST1--CST3,
les relations $(\mathrm{MO}_{\mathrm{II}})$, $(\mathrm{MO}_{\mathrm{III}})$, et
la \emph{transportabilité} de la relation typique) sont fournis comme
\textbf{hypothèses honnêtes explicites}, jamais postulés en axiomes
--- \texttt{theorie\_ensembles()==22} reste vrai. Le statut dominant
« conditionnel honnête » dit donc : \emph{si les schémas de transport valent,
alors tel critère en découle}. Le c\oe ur $\sigma/\Sigma$ (récurrence sur le
schéma d'échelon) et l'existence effective des structures transportées sont
signalés comme résidus (\S\,Statut).

\subsection{IV.1 --- Espèces, isomorphismes, transport de structure}

\textbf{Transport et isomorphisme.} Si $f=(f_1,\dots,f_n)$ est une famille de
bijections $E_i\!\to\!E_i'$, alors $f$ est un isomorphisme de $(E,U)$ sur
$(E,U')$ où $U'=\langle f\rangle^{S}(U)$ est la structure transportée
(\texttt{transport\_donne\_isomorphisme}, c\oe ur objet de CST5). La trivialité
de réflexivité $U'=U'$ de la clause $(4)$ est, elle, \emph{close}
(\texttt{transport\_egalite}).

\textbf{Groupoïde des isomorphismes.} Les propriétés de groupoïde sont
certifiées (toutes \emph{conditionnelles honnêtes}, modules
\texttt{ensembles\_transport\_iso\_props.py} et \texttt{ensembles\_structures\_props.py}) :
\begin{itemize}
  \item \emph{réflexivité} : $\langle\Delta_E\rangle^{S}(U)=U \vdash$
        $(E,U)$ est isomorphe à lui-même, et $\Delta_E$ est un automorphisme ;
  \item \emph{réciproque} (CST3) : $\bigl(\langle f\rangle^{S}\bigr)^{-1}=
        \langle f^{-1}\rangle^{S}$, d'où « la réciproque d'un isomorphisme est un
        isomorphisme » ;
  \item \emph{composition} (CST4) : la composée de deux isomorphismes est un
        isomorphisme ($\langle g\circ f\rangle^{S}=\langle g\rangle^{S}\circ
        \langle f\rangle^{S}$ en prémisse) ;
  \item \emph{unicité du transport} (CST5) : deux structures transportées par le
        même isomorphisme coïncident ($V=V'$).
\end{itemize}

\textbf{Déduction d'espèces (IV.1.6--7).} La fonctorialité de la déduction
(CST6 : un isomorphisme se transporte en isomorphisme de l'espèce déduite) et
l'équivalence des isomorphismes pour des espèces équivalentes (CST7) sont
certifiées conditionnellement.

\subsection{IV.2 --- Morphismes et structures dérivées}

La section la plus fournie. Le préordre « \emph{plus fine} » sur les structures
de même base est certifié \textbf{réflexif} et \textbf{transitif} (IV.2.2,
d'après $(\mathrm{MO}_{\mathrm{III}})$ et $(\mathrm{MO}_{\mathrm{II}})$), la
transitivité pleine étant obtenue après normalisation $\Delta_E\circ\Delta_E=
\Delta_E$. La composée de deux $\sigma$-morphismes est un $\sigma$-morphisme
($(\mathrm{MO}_{\mathrm{II}})$).

\textbf{Structures universelles (initiales, finales).} Les sens « faciles » de
$(\mathrm{IN})$ et $(\mathrm{FI})$ sont prouvés (la structure initiale rend les
$f_\iota$ des morphismes ; dual pour la finale), puis l'\textbf{unicité} :
\begin{itemize}
  \item deux structures initiales sont \emph{mutuellement plus fines}, d'où
        $\mathcal{I}=\mathcal{I}'$ sous antisymétrie (\textbf{CST9}) ; dual
        $\mathcal{F}=\mathcal{F}'$ (\textbf{CST18}) ;
  \item transitivité des structures \emph{initiales} (CST10), \emph{induites}
        (CST11), \emph{finales} (CST19), \emph{quotient} (CST21) --- chacune
        ramenée à un palier d'égalité ;
  \item \textbf{CST12} (restriction d'un morphisme aux sous-structures) et
        \textbf{CST20} (passage des morphismes aux quotients), duaux l'un de
        l'autre.
\end{itemize}

\textbf{Produits et graphes.} La structure produit : associativité (CST13),
compatibilité produit/sous-structure (CST14), image réciproque du produit
(CST15) ; une \emph{famille} de morphismes définit un morphisme dans le produit
(\textbf{CST16}) ; et un morphisme est caractérisé par son graphe
(\textbf{CST17}). Enfin la \emph{décomposition canonique} d'une application
(IV.2.5) fournit des lemmes d'extraction \textbf{clos}
(\texttt{ensembles\_structures\_residus.py}).

\subsection{IV.3 --- Applications universelles}

L'équivalence définissant une application universelle est \textbf{close} :
\[
  (\mathrm{AU})\ \Longleftrightarrow\ (\mathrm{AU}_{\mathrm I}')\wedge
  (\mathrm{AU}_{\mathrm{II}}')
\]
(existence d'une factorisation et unicité ; modules
\texttt{ensembles\_universel\_applications.py}, trois théorèmes clos). L'unicité
de la solution universelle \emph{à un unique isomorphisme près} (CST8) est
certifiée conditionnellement, et les extractions logiques de l'hypothèse de
CST22 ($\mathrm{CU}_{\mathrm I}$, $\mathrm{CU}_{\mathrm{III}}$) sont closes
(\texttt{ensembles\_structures\_complements.py}, E~IV.16).

\subsection{Statut et résidus honnêtes}

Sont \textbf{certifiés clos} : la trivialité de transport, l'équivalence
$(\mathrm{AU})\Leftrightarrow(\mathrm{AU}_{\mathrm I}')\wedge(\mathrm{AU}_{\mathrm{II}}')$,
et les lemmes d'extraction (décomposition canonique, CU de CST22). Sont
\textbf{conditionnels honnêtes} la quasi-totalité des critères CST, les schémas
de transport étant prémisses explicites. Sont \textbf{reportés}, et présentés
comme tels :
\begin{itemize}
  \item la \emph{transportabilité effective} de la relation typique $R$ --- donc
        l'\emph{existence} même des structures transportées, initiales, finales,
        induites et quotient (\textbf{CST22}) ;
  \item les preuves \emph{par récurrence sur le schéma d'échelon} (fonctorialité
        effective de l'extension d'échelon, CST1--CST3 dans leur version
        constructive) ;
  \item \textbf{CST23} dans sa forme pleine (le lemme de contraposition
        ponctuelle présent n'en est \emph{pas} le contenu, seulement le squelette
        logique).
\end{itemize}

\textbf{Piège neutralisé.} Un lemme $\langle\,\cdot\,\rangle$ se réduisant à une
tautologie $P\Leftrightarrow P$
(\texttt{\_est\_iso\_morph\_reflexivite\_triviale}) a été identifié et maintenu
\emph{hors} de l'interface publique (\texttt{\_\_all\_\_}) : il n'est jamais
présenté comme le critère $(\mathrm{MO}_{\mathrm{III}})$, conformément à la règle
« aucune tautologie déguisée en théorème » (cf.\ annexe, Journal des erreurs).

\subsection{§IV.2 --- CST8 : un morphisme inversible est un isomorphisme}
\textbf{Énoncé (Bourbaki, E~IV.12, §2.1).} « Soit $f$ un $\sigma$-morphisme de $E$ dans $E'$,
$g$ un $\sigma$-morphisme de $E'$ dans $E$. Si $g\circ f=\mathrm{Id}_E$ et $f\circ g=\mathrm{Id}_{E'}$,
$f$ est un isomorphisme de $E$ sur $E'$ et $g$ l'isomorphisme réciproque. »

\textbf{Formalisé.} $\{\,\mathrm{morph}(E,\mathscr S,E',\mathscr S',f),\
\mathrm{morph}(E',\mathscr S',E,\mathscr S,g),\ g=f^{-1}\,\}\vdash
\mathrm{est\_iso\_morph}(E,\mathscr S,E',\mathscr S',f)$. L'inversibilité bilatère
$g\circ f=\mathrm{Id}$, $f\circ g=\mathrm{Id}$ est résumée fidèlement par son conséquent
$g=f^{-1}$ (corollaire II, p.\,18, brique reportée comme CST3/12/20). Module :
\texttt{.../iv\_2\_morphismes\_structures\_derivees/ensembles\_cst8\_inversible\_iso.py}.

\textbf{Preuve (\Nstar).} $S6$/Leibniz réécrit $\mathrm{morph}(\dots,g)$ en
$\mathrm{morph}(\dots,f^{-1})$ ; conjonction avec $\mathrm{morph}(\dots,f)$. \emph{Fidélité} :
le \emph{vrai} CST8 (E~IV.12, §2.1) était \emph{absent} --- les fonctions code nommées « CST8 »
encodaient en réalité l'unicité de la solution universelle (§IV.3.1, E~IV.23) ; l'étiquette
\texttt{@livre} fautive a été corrigée.

\textbf{Statut.} CLOS sous $3$ hypothèses ; \texttt{theorie\_ensembles()==22}.
